\section{String Theory from Observer Patch Holography}\label{string-theory-from-observer-patch-holography}

\begin{quote}
This part studies a string/worldsheet reorganization of the OPH edge-sector branch. Its input is the fixed-cutoff collar package, the heat-kernel sector weights, the conditional gravity branch, and, where needed, the compact-gauge branch. The additional premises are the stated large-\(N_{\mathrm{edge}}\) and modular-sewing assumptions.
\end{quote}

\subsection{Scope and Input}\label{a-complete-derivation}

The bosonic edge-sector branch admits a string-theoretic description of its effective degrees of freedom if the stated heat-kernel, large-\(N\), and modular-sewing premises are realized.

\begin{center}\rule{0.5\linewidth}{0.5pt}\end{center}

\section{Overview of the Route}\label{1-overview-what-we-derive}

The chain studied in this section is:

\begin{verbatim}
Five OPH axioms + regulator premises used in the edge-sector branch
    \ensuremath{\downarrow}
Edge-Center Decomposition + Heat-Kernel Sector Weights
    \ensuremath{\downarrow}
2D Yang-Mills Theory on the Holographic Screen
    \ensuremath{\downarrow}
String Theory (Worldsheet Expansion via Gross-Taylor)
\end{verbatim}

OPH also supplies:

\begin{itemize}
\tightlist
\item
  \textbf{Worldsheet kinematics}: Conformal/modular sewing from geometric modular flow
\item
  \textbf{Target-space dynamics}: conditional scaling-limit Einstein branch from entanglement equilibrium
\item
  \textbf{Higher gauge structure}: The natural slot for B-field/gerbe data
\end{itemize}

This gives a conditional route by which the OPH edge/boundary sector admits a string-theoretic reorganization if the stated large-\(N\) and modular-sewing premises are realized.

\begin{center}\rule{0.5\linewidth}{0.5pt}\end{center}

\section{OPH Input}\label{2-oph-inputs-used}

The analysis uses four ingredients from the main OPH trunk: the fixed-cutoff edge-center package, the heat-kernel edge-sector weights, the conditional geometric-modular gravity branch, and the realized compact-gauge branch whenever gauge content is mentioned. The additional inputs specific to this section are the stated large-\(N_{\mathrm{edge}}\) and modular-sewing premises.

\begin{center}\rule{0.5\linewidth}{0.5pt}\end{center}

\section{Target-Space Dynamics}\label{3-einstein-equation-from-entanglement-equilibrium}

On the \(T2/T3\) branch together with the theorem-local geometric cap-pair extraction and ordered cut-pair rigidity inputs, and with the stated downstream null-stress identification, OPH yields the target-space equation
\[
\boxed{G_{ab} + \Lambda g_{ab} = 8\pi G \langle T_{ab} \rangle}.
\]
This equation supplies the target-space dynamics used in the worldsheet analysis. The additional question is whether the edge-sector and large-\(N_{\mathrm{edge}}\) data admit a controlled string reorganization compatible with that target-space branch.

\begin{center}\rule{0.5\linewidth}{0.5pt}\end{center}

\section{Edge Sectors and 2D Yang-Mills}\label{4-edge-sectors-and-2d-yang-mills}

\subsection{Edge-Sector Structure}\label{41-edge-sector-structure}

From the OPH edge-center decomposition, the collar Hilbert space has structure: \[\mathcal{H}_{\text{collar}} \cong \bigoplus_{\alpha} \left(\mathcal{H}_{b_L^\alpha} \otimes \mathcal{H}_{b_R^\alpha}\right)\]

with \(\alpha\) labeling irreps/charges of the boundary gauge group.

This is exactly the structure of gauge theory with boundary: cutting gauge space produces edge modes whose labels are representation data.

\subsection{Heat-Kernel Weights from MaxEnt}\label{42-heat-kernel-weights-from-maxent}

OPH\textquotesingle s MaxEnt + bi-invariant constraints give sector probabilities: \[p_R(t) \propto d_R \, e^{-t \, C_2(R)}\]

Because of the Peter-Weyl doubling (loops see both factors), the effective weight is: \[w_R(t) \propto d_R \, p_R(t) \propto d_R^2 \, e^{-t \, C_2(R)}\]

\subsection{Edge Partition Function}\label{43-edge-partition-function}

Define the edge partition function at diffusion parameter \(t\): \[Z_{\text{edge}}(t) := \sum_R d_R^2 \, e^{-t \, C_2(R)}\]

\subsection{Identification with Group Heat Kernel}\label{44-identification-with-group-heat-kernel}

The heat kernel on compact group \(G\), expanded in characters: \[K_t(U) = \sum_R d_R \, \chi_R(U) \, e^{-t C_2(R)}\]

At the identity \(U = 1\) (using \(\chi_R(1) = d_R\)): \[K_t(1) = \sum_R d_R^2 \, e^{-t C_2(R)} = Z_{\text{edge}}(t)\]

\textbf{Result:} The OPH edge ensemble is literally the heat kernel at the identity.

\subsection{Gluing = Markov = Area Additivity}\label{45-gluing--markov--area-additivity}

Heat kernels satisfy the Chapman-Kolmogorov (gluing) property: \[K_{t_1 + t_2}(U) = \int_G dV \, K_{t_1}(UV^{-1}) \, K_{t_2}(V)\]

This matches OPH patch sewing: add diffusion parameters when gluing collars/strips.

\subsection{Exact Match with 2D Yang-Mills}\label{46-exact-match-with-2d-yang-mills}

The partition function of 2D Yang-Mills on a closed genus-\(g\) surface \(\Sigma_g\) with area \(A\): \[Z_{\text{2D YM}}(\Sigma_g) = \sum_R (d_R)^{2-2g} \exp\left[-\frac{g_{\text{YM}}^2 A}{2} C_2(R)\right]\]

For \(\Sigma_0 = S^2\) (genus 0): \[Z_{\text{2D YM}}(S^2) = \sum_R d_R^2 \exp\left[-\frac{g_{\text{YM}}^2 A}{2} C_2(R)\right]\]

\textbf{Identification:} \[t = \frac{g_{\text{YM}}^2 A}{2}\]

\textbf{Theorem 4.1:} OPH edge-sector MaxEnt gives the 2D Yang-Mills heat-kernel partition-function form for the edge branch, with collar diffusion parameter \(t\) playing the role of "coupling \ensuremath{\times} area". In the effective 2D Yang-Mills regime, this is the standard 2D YM identification.

\begin{center}\rule{0.5\linewidth}{0.5pt}\end{center}

\section{2D Yang-Mills as String Theory}\label{5-2d-yang-mills-as-string-theory}

\subsection{The Gross-Taylor Expansion}\label{51-the-gross-taylor-expansion}

In the large-\(N\) limit for \(G = SU(N)\), 2D Yang-Mills admits an exact reorganization as a sum over worldsheets~\cite{grossTaylor93}.

The partition function expands as: \[\log Z = \sum_{h \geq 0} N^{2-2h} F_h(\lambda), \quad \lambda := g_{\text{YM}}^2 N\]

where:

\begin{itemize}
\tightlist
\item
  \(h\) is the worldsheet genus
\item
  \(g_s \sim 1/N\) is the string coupling
\item
  Contributions scale as \(N^{2-2h}\) for worldsheet genus \(h\)
\end{itemize}

\subsection{Worldsheet Interpretation}\label{52-worldsheet-interpretation}

The \(1/N\) expansion coefficients count \textbf{branched coverings} of the target surface,i.e., "string configurations without folds."

On that conditional large-\(N\) branch, this is a genuine string-theoretic reorganization:

\begin{itemize}
\tightlist
\item
  \textbf{Target space:} The 2D screen surface
\item
  \textbf{Worldsheet:} The covering surface summed over
\item
  \textbf{String coupling:} \(g_s \sim 1/N\)
\end{itemize}

\subsection{\texorpdfstring{The Large-\(N\) Regime}{The Large-N Regime}}\label{53-the-large-n-regime}

The Gross-Taylor expansion is controlled in the large-\(N\) limit. OPH provides an effective large-\(N\) parameter from screen microstructure: many independent edge channels/pixels whose combined algebra behaves like large matrix degrees of freedom. The effective edge Hilbert space dimension scales as: \[\dim(\tilde{\mathcal{H}}_\partial) \sim N^{\#(\text{boundary cells})}\]

so \(1/N\) becomes the handle-counting coupling.

\subsection{Genus Expansion}\label{54-genus-expansion}

Define generating functional: \[F(t) := \log Z_{\text{edge}}(t)\]

In controlled large-\(N\) regime: \[F(t) = \sum_{h \geq 0} N^{2-2h} F_h(\lambda), \quad \lambda := g_{\text{YM}}^2 N\]

This is therefore a perturbative closed-string genus expansion on the same conditional branch.

\begin{center}\rule{0.5\linewidth}{0.5pt}\end{center}

\section{Worldsheet Kinematics from Modular Flow}\label{6-worldsheet-kinematics-from-modular-flow}

\subsection{Conformal Structure}\label{61-conformal-structure}

OPH provides the conformal/Lorentz data in two logically distinct steps.

\textbf{Screen conformal group.}
\[
\mathrm{Conf}^+(S^2)\cong \mathrm{PSL}(2,\mathbb C)\cong \mathrm{SO}^+(3,1).
\]
This is the ambient conformal group of the celestial sphere, interpreting screen angles as null directions of \(3+1\)D spacetime.

\textbf{Geometric modular branch.} Lorentz kinematics uses that group only on the explicit \textbf{T1}+\textbf{T3} branch where the emitted scaling-limit cap pair has the standard geometric modular automorphism group with the \(2\pi\) normalization. OPH does not derive that branch from MaxEnt alone, and it does not claim that the refinement-limit cap algebra remains type I. In the generic non-type-I case the relevant geometric modular action is outer rather than an inner density-matrix formula, so any stronger vacuum/canonical identification remains theorem-external.
\subsection{Sewing Structure}\label{62-sewing-structure}

String worldsheets are sewn along boundaries/cuts. Consistent sewing requires conformal/modular data.

OPH\textquotesingle s Markov/additivity theorems on collars derive this exact factorization structure that makes sewing well-behaved.

\subsection{The Analogy with Standard String Theory}\label{63-the-analogy-with-standard-string-theory}

{\def\LTcaptype{none} % do not increment counter
\begin{longtable}[]{@{}>{\RaggedRight\arraybackslash}p{0.480\textwidth}>{\RaggedRight\arraybackslash}p{0.480\textwidth}@{}}
\toprule\noalign{}
String Theory & OPH \\
\midrule\noalign{}
\endhead
\bottomrule\noalign{}
\endlastfoot
Worldsheet conformal symmetry & Screen conformal group \\
Modular invariance & Overlap consistency \\
Target-space Einstein eq. & Entanglement equilibrium \\
Sewing axioms & Markov collar gluing \\
\end{longtable}
}

\begin{center}\rule{0.5\linewidth}{0.5pt}\end{center}

\section{The B-Field Slot: Higher Gauge Structure}\label{7-the-b-field-slot-higher-gauge-structure}

\subsection{Nonabelian Gluing Cocycles}\label{71-nonabelian-gluing-cocycles}

OPH includes gluing by a crossed-module 2-cocycle \((g_{ij}, h_{ijk})\) on triple overlaps, with coherence on quadruple overlaps.

This is precisely the higher-gauge/gerbe structure associated in string theory with:

\begin{itemize}
\tightlist
\item
  2-form gauge field \(B\) (Kalb-Ramond field)
\item
  Gauge invariance \(B \to B + d\eta\)
\end{itemize}

\subsection{Strictification and Trivialization}\label{72-strictification-and-trivialization}

OPH shows strictification exists iff the abelian \v{C}ech \(H^2\) class is trivial.

This parallels: "background \(B\)-field is a gerbe class; strings couple to it; trivialization conditions matter."

\subsection{D-Brane Interpretation}\label{73-d-brane-interpretation}

Certain defects/boundaries where the class trivializes are the seed of D-brane-like boundary conditions.

The mathematical structure where string theory places B-field/gerbe data is directly present in OPH. The full open/closed string boundary state formalism follows from extending this framework to include D-brane-like boundary conditions.

\begin{center}\rule{0.5\linewidth}{0.5pt}\end{center}

\section{OPH String Scale and Parameters}\label{8-oph-string-scale-and-parameters}

\subsection{String Scale from Pixel Area}\label{81-string-scale-from-pixel-area}

OPH has a UV length from pixel area: \[a_{\text{cell}} \approx 1.63 \, \ell_P^2, \quad \ell_{\text{UV}} = \sqrt{a_{\text{cell}}} \approx 1.28 \, \ell_P\]

Natural identification: \[\alpha' \sim \ell_{\text{UV}}^2 \sim a_{\text{cell}}\]

String tension: \[T_{\text{string}} = \frac{1}{2\pi\alpha'} \sim \frac{1}{2\pi a_{\text{cell}}}\]

Numerically (in Planck units): \[T_{\text{string}} \approx \frac{1}{2\pi \cdot 1.63} \, \ell_P^{-2} \approx 0.098 \, \ell_P^{-2}\]

This is order-Planck tension, as expected if UV cutoff sits near \(\ell_P\).

\subsection{Worldsheet YM Coupling}\label{82-worldsheet-ym-coupling}

From OPH: \[g_{\text{ent}}^2 = \frac{t}{2\pi}\]

In 2D YM, heat-kernel time is \(t \sim (g_{\text{YM}}^2 A)/2\).

So OPH\textquotesingle s \(t\) really does play the role of 2D YM area-time / diffusion parameter.

\subsection{Current Algebra Central Charge (WZW)}\label{83-current-algebra-central-charge-wzw}

If edge gauge group embeds into worldsheet current algebra (WZW-type): \[c = \frac{k \, \dim\mathfrak{g}}{k + h^\vee}\]

For \(SU(3) \times SU(2) \times U(1)\) with \(h^\vee_{SU(3)} = 3\), \(h^\vee_{SU(2)} = 2\): \[c_{\text{int}}(k) = \frac{8k}{k+3} + \frac{3k}{k+2} + 1\]

For 4D superstring compactification with internal \(c_{\text{int}} \approx 9\): \[k = \frac{5 + \sqrt{89}}{2} \approx 7.22\]

Integer levels \(k = 7\) or \(8\) bracket \(c_{\text{int}} \approx 9\).

\textbf{Conclusion:} OPH gauge content is numerically consistent with reasonable-level current algebra sector inside critical worldsheet theory.

\begin{center}\rule{0.5\linewidth}{0.5pt}\end{center}

\section{From 2D Yang-Mills to Critical Superstrings}\label{9-from-2d-yang-mills-to-critical-superstrings}

The derivation chain from OPH to string theory proceeds in two stages.

\subsection{Established Results}\label{91-established-results}

The following are derived from the five OPH axioms together with the stated regulator, heat-kernel, and large-\(N\) premises:

\begin{enumerate}
\def\labelenumi{\arabic{enumi}.}
\tightlist
\item
  Finite microscopic DoF on the screen, type-I local algebras, boundary gauge invariance
\item
  EC decomposition and Markov structure on collars
\item
  Sector category and reconstruction of the compact gauge group, with \(SU(3) \times SU(2) \times U(1)/\mathbb{Z}_6\) selected under the MAR admissibility package on the realized gauge branch
\item
  Heat-kernel edge-sector weights \(p_R \propto d_R e^{-t C_2(R)}\)
\item
  Identification with the 2D Yang-Mills partition function (Theorem 4.1)
\item
  Conditional geometric modular flow with \(2\pi\) normalization on the extracted prime geometric subnet under \textbf{T2} plus cap-pair extraction and ordered cut-pair rigidity, and Lorentz symmetry from Conf\((S^2)\) on that branch
\item
  Conditional scaling-limit Einstein branch from entanglement equilibrium
\item
  Nonabelian gluing (crossed module/2-group) providing higher gauge data
\end{enumerate}

The genus expansion \(\log Z = \sum_{h \geq 0} N^{2-2h} F_h(\lambda)\) follows from the Gross-Taylor rewriting~\cite{grossTaylor93} of 2D Yang-Mills in the large-\(N\) regime, with \(N\) identified as the effective edge Hilbert space parameter.

\subsection{Extensions to Critical Superstrings}\label{92-extensions-to-critical-superstrings}

The full path from OPH to critical superstring theory requires:

\begin{itemize}
\tightlist
\item
  \textbf{Worldsheet CFT}: Virasoro emergence from modular data, modular invariance from overlap consistency
\item
  \textbf{Complete massless spectrum}: metric + B-field + dilaton + RR sectors, anomaly cancellation from worldsheet beta functions
\item
  \textbf{Internal sector matching}: the SM gauge content with \(c_{\text{int}} \approx 9\) is numerically consistent with integer WZW levels \(k = 7\) or \(8\) (\S{}8.3)
\end{itemize}

The graviton is contained implicitly via the Einstein equation derived from entanglement equilibrium.

\begin{center}\rule{0.5\linewidth}{0.5pt}\end{center}

\section{Summary}\label{10-summary}

The string/worldsheet material belongs to Phase III. The recovered-core claim set of OPH is Phase
I, namely D1--D5 together with D7--D9, while the present section requires the additional
large-\(N_{\mathrm{edge}}\), modular-sewing, and worldsheet-CFT premises recorded below.

The theorem boundary is therefore:

\begin{enumerate}
\def\labelenumi{\arabic{enumi}.}
\item
  \textbf{Secured upstream input.} The theorem-level statement actually available before any string reorganization is that the edge-sector weights have the 2D Yang-Mills heat-kernel form. That identity is the robust bridge from OPH edge data to a worldsheet-like reorganization.
\item
  \textbf{Conditional continuation.} If a controlled large-\(N_{\mathrm{edge}}\) regime exists, with \(N_{\mathrm{edge}}\) distinct from the physical color rank \(N_c=3\), then the Gross-Taylor rewriting~\cite{grossTaylor93} turns the edge partition function into a genus expansion. This is the precise theorem boundary of the worldsheet/string reorganization claim.
\item
  \textbf{Not yet claimed.} Critical superstring worldsheet CFT, full massless spectrum, anomaly cancellation, and complete internal-sector matching all require further premises beyond the recovered OPH core.
\end{enumerate}

The broader OPH framework supplies suggestive ingredients for such a continuation, including
modular geometry, a conditional Einstein branch, and higher-gauge/crossed-module data. Those
ingredients do not by themselves upgrade this part beyond Phase III. Failure of the branch would
retract only this continuation and would not alter the relativity-plus-Standard-Model recovery
package.

\begin{center}\rule{0.5\linewidth}{0.5pt}\end{center}

\section{Appendix A: Entanglement First Law Numerical Check}\label{appendix-a-entanglement-first-law-numerical-check}

The "entanglement first law" step (\(\delta S = \delta\langle K \rangle\)) is load-bearing in the Einstein derivation. Here we verify it numerically.

\subsection{A.1 Setup}\label{a1-setup}

Generate random bipartite density matrix on \(4 \times 4\) Hilbert space (2 qubits \ensuremath{\times} 2 qubits). Perturb by small \(\varepsilon\) and compare:

\[\Delta S_A := S(\rho_A(\varepsilon)) - S(\rho_A(0))\] \[\Delta\langle K_A \rangle := \text{Tr}[(\rho_A(\varepsilon) - \rho_A(0)) K_A]\]

where \(K_A = -\log \rho_A(0)\).

\subsection{A.2 Expected Result}\label{a2-expected-result}

As \(\varepsilon \to 0\): \[\frac{|\Delta S_A - \Delta\langle K_A \rangle|}{|\Delta\langle K_A \rangle|} \to 0\]

linearly in \(\varepsilon\) (first-order identity).

\subsection{A.3 Code}\label{a3-code}

See \texttt{code/entanglement\_first\_law.py} for implementation.

\begin{center}\rule{0.5\linewidth}{0.5pt}\end{center}

\section{Appendix B: Key Equations Summary}\label{appendix-b-key-equations-summary}

{\def\LTcaptype{none} % do not increment counter
\begin{longtable}[]{@{}>{\RaggedRight\arraybackslash}p{0.480\textwidth}>{\RaggedRight\arraybackslash}p{0.480\textwidth}@{}}
\toprule\noalign{}
Result & Equation \\
\midrule\noalign{}
\endhead
\bottomrule\noalign{}
\endlastfoot
Generalized entropy & \(S_{\text{gen}} = \frac{A}{4G} + S_{\text{bulk}}\) \\
Einstein from equilibrium & \(G_{ab} + \Lambda g_{ab} = 8\pi G \langle T_{ab} \rangle\) \\
Edge-sector weights & \(p_R(t) \propto d_R \, e^{-t C_2(R)}\) \\
Edge partition function & \(Z_{\text{edge}}(t) = \sum_R d_R^2 \, e^{-t C_2(R)} = K_t(1)\) \\
2D YM partition function & \(Z_{\text{2D YM}}(S^2) = \sum_R d_R^2 \, e^{-\frac{g^2 A}{2} C_2(R)}\) \\
Parameter identification & \(t = \frac{g_{\text{YM}}^2 A}{2}\) \\
String tension & \(T_s = \frac{1}{2\pi\alpha'} \sim \frac{1}{2\pi a_{\text{cell}}}\) \\
Genus expansion & \(\log Z = \sum_{h \geq 0} N^{2-2h} F_h(\lambda)\) \\
\end{longtable}
}

\begin{center}\rule{0.5\linewidth}{0.5pt}\end{center}

\section{References}\label{references-string-theory}

This derivation uses:

\begin{enumerate}
\def\labelenumi{\arabic{enumi}.}
\tightlist
\item
  Canonical OPH axiom package and theorem-local regulator premises from the main OPH trunk
\item
  Heat-kernel / edge-sector results from the fixed-cutoff collar and gauge branch
\item
  Gross-Taylor 2D YM / string duality~\cite{grossTaylor93}
\item
  Jacobson entanglement equilibrium~\cite{jacobson95,jacobson16}
\item
  Standard results: Tomita-Takesaki modular theory, Peter-Weyl theorem, Raychaudhuri equation
\end{enumerate}
